{{دیوان
| عنوان = اِبلیس کا فرمان اپنے سیاسی فرزندوں کے نام
| ماخذ = ویکی ماخذ
| ماخذ_لنک = https://ur.wikisource.org/wiki/%D8%A7%D9%90%D8%A8%D9%84%DB%8C%D8%B3_%DA%A9%D8%A7_%D9%81%D8%B1%D9%85%D8%A7%D9%86_%D8%A7%D9%BE%D9%86%DB%92_%D8%B3%DB%8C%D8%A7%D8%B3%DB%8C_%D9%81%D8%B1%D8%B2%D9%86%D8%AF%D9%88%DA%BA_%DA%A9%DB%92_%D9%86%D8%A7%D9%85
| مقابلہ = اقبال سائبر لائبریری (اقبال اکادمی پاکستان)
| مقابلہ_لنک = https://www.iqbalcyberlibrary.net
}}
<poem>
لا کر بَرہمنوں کو سیاست کے پیچ میں
زُناّریوں کو دَیرِ کُہن سے نکال دو

وہ فاقہ کش کہ موت سے ڈرتا نہیں ذرا
رُوحِ محمدؐ اس کے بدن سے نکال دو

فکرِ عرب کو دے کے فرنگی تخیّلات
اسلام کو حِجاز و یمن سے نکال دو

افغانیوں کی غیرتِ دیں کا ہے یہ علاج
مُلّا کو اُن کے کوہ و دمن سے نکال دو

اہلِ حرم سے اُن کی روایات چھِین لو
آہُو کو مرغزارِ خُتن سے نکال دو

اقبالؔ کے نفَس سے ہے لالے کی آگ تیز
ایسے غزل سرا کو چمن سے نکال دو!
</poem>

٭ بھوپال (شیش محل) میں لکھے گئے
